\documentclass[12pt]{article}
\usepackage[utf8]{inputenc}
\usepackage[T1]{fontenc}
\usepackage{hyperref}
\title{Lydian Language Resources for K-12 Curriculum Development}
\author{Compiled by Claude Sonnet 5.5 (claude-sonnet-5-5) for Andr\'as Kornai}
\date{}
\begin{document}
\maketitle

\section{Introduction}
Lydian (ISO 639-3: xld; Glottocode: lydi1241) is an extinct Anatolian Indo-European language of ancient Lydia in western Anatolia. Most of its texts come from Sardis and date mostly to the 5th and 4th centuries BCE. They are written in an alphabet related to the Greek one. Fewer than 30 texts are long or fairly complete. Enno Littmann deciphered the language in 1916 with the help of the Aramaic--Lydian bilingual from Sardis. The 109 inscriptions known by 1986 were treated comprehensively by Gusmani.

\section{Primary Educational Resources}
I found no school curricula, primers or textbooks for Lydian.

\section{Linguistic and Reference Materials}
\begin{itemize}
\item Lydian language (Wikipedia). \url{https://www.Wikipedia.org/wiki/Maeonian}
\item eDiAna Lydian dictionary entry. \url{https://www.ediana-remake.gwi.uni-muenchen.de/dictionary/471}
\item A second eDiAna Lydian dictionary entry. \url{https://www.ediana-remake.gwi.uni-muenchen.de/dictionary/473}
\item eDiAna Lydian corpus entry. \url{https://www.ediana-remake.gwi.uni-muenchen.de/corpus/lydian/66}
\item Britannica, Anatolian languages using an alphabet. \url{https://britannica.com/topic/Anatolian-languages/Languages-using-an-alphabet}
\item Academia.edu item from the search, probably Yakubovich's work on the Sardis stele LW 22. I did not open it, so check it before relying on it. \url{https://www.academia.edu/35571454}
\item Overview page ``Lydian: The Anatolian Branch of Indo-European and Its Documentation''. It is a low-authority source and should be checked. \url{https://gravitee.io/corpus/gen-1871/anatolian-languages/lydian.html}
\end{itemize}

\section{Religious and Community Resources}
There is no living community. The surviving texts are mainly tomb inscriptions, with some decrees and verse pieces.

\section{Digital and Audio Resources}
eDiAna provides online dictionary and corpus entries. A Carlsberg Foundation-funded project at Leiden, ``Lost Voices of Lydia'' (led by Tobias Søborg), is preparing a new edition of all known Lydian inscriptions, with a dictionary and a grammar. \url{https://carlsbergfondet.dk/en/what-we-have-funded/cf24-0951/} I found no audio resources.

\section{Curriculum Development Context}
Lydian is a classical-language case with a small and partly opaque corpus. A heritage unit could cover the script, Sardis and the Lydian kingdom, and the decipherment story. The new Leiden edition may later provide better reference material. I found no existing school materials.

\end{document}